\question[3]
Die Hilfe zu \texttt{print} beginnt so:
\begin{lstlisting}[language={}]
print(*args, sep=' ', end='\n', file=None, flush=False)
    Prints the values to a stream, or to sys.stdout by default.
    sep
      string inserted between values, default a space.
    end
      string appended after the last value, default a newline.
\end{lstlisting}
\begin{parts}
\part Was erscheint auf der Konsole?
\begin{lstlisting}
print(1, 2, 3, sep='-', end='!')
print('x')
\end{lstlisting}
\part Schreibe einen Aufruf von \texttt{print} mit den drei Argumenten \texttt{'a'}, \texttt{'b'}, \texttt{'c'}, der \texttt{a:b:c} ausgibt.
\end{parts}
\begin{solutionbox}{2cm}
\begin{lstlisting}[language={}]
a) 1-2-3!x
b) print('a', 'b', 'c', sep=':')
\end{lstlisting}
\end{solutionbox}
